\section{Trade-Off Analysis}
\subsection{Concept Trade-Off}
 
\begin{table}[H]
\centering
\begin{tabular}{|c|c|c|c|}
\hline
\textbf{Attribute} & \textbf{ORACLE} & \textbf{EVA} & \textbf{NOVAE} \\ \hline
 Orbit Type      &      LEO     &     HEO      &     L2      \\ \hline
 $m_{fuel}$ required (kg)*     &    350,000      &      400,000     &   900,000         \\ \hline
Energy Source       &      Solar Arrays    &      Solar Arrays     &    Solar Arrays       \\ \hline
 Solar Panels Area (m$^2$)      &    $10-20$       &    $5-15$       &    $20-30$       \\ \hline
   Estimated Dry Mass (kg)       &       $\leq$ 6,200    &   $\leq$ 6,200        &        $\leq$ 11,100   \\ \hline
    Antenna Gain      &       High    &    High       &      High     \\ \hline
    Antenna Bit Rate      &     Medium      &     Medium      &     High      \\ \hline
    \# of User Needs Met      &     10      &      10     &     9.5      \\ \hline
\end{tabular}
\caption{Attribute's summary for the three concepts proposed}
\end{table}
\textbf{*}: Fuel is required to place the satellite into a parking orbit. We estimate this based on \textbf{Merlin 1D engine's} average fuel burn per kg.

\subsection{Launch Site and Vehicle}
Table (\ref{LaunchVehicle}) lists down the most openly available launch systems in the market as of November 2024; this is based on political relations hence launch systems by China and Russia are exempted. 
\begin{table}[h]
\centering
\begin{tabular}{|c|c|c|c|c|c|}
\hline
\textbf{Name} & \textbf{Company} & \textbf{Pld.(kg)} & \textbf{Cost (\$/kg)} & \textbf{Sustainability} & \textbf{Success Rate} \\ \hline
Falcon Heavy & SpaceX & 63,000 & 1,400 & High& 100\% \\ \hline
Falcon 9 & SpaceX & 22,800 & 1,520 & High& 98.6\% \\ \hline
S.L.S. & NASA & 130,000 & 43,157 &Low & 100\% \\ \hline
Arianne A62 & ESA & 10,350 & ($-$) & Medium& ($-$)\\ \hline
Arianne A64 & ESA & 21,650 & ($-$) & Medium& ($-$)\\ \hline
Atlas V& NASA & 18,850 & 6,000 & Low& 100\% \\ \hline
Starship & SpaceX & 150,000 & 20 & High& ($-$) \\ \hline
Vulcan Centaur & ULA & 27,200 &5,514 & Medium&($-$) \\ \hline
\end{tabular}
\caption{Summary table for the launch vehicle architecture. Used for decision-making. Pld. is max payload. }
\label{LaunchVehicle}
\end{table}

Taking into account the predicted weight of our satellites and the more favourable prices offered by SpaceX, the Falcon 9 shall be the preferred option for our launch system, with the launch site in the Kennedy Space Center in Cape Canaveral, USA. Given the historic decrease in the cost over the years, it is highly likely that over the 5 years, these will further decrease. \\
\vspace{2mm}
The trade-off between the three proposed concepts is done based on the following parameters:
\begin{itemize}
    \item \textbf{Cost}: all expenses associated with the design, development, launch, operation, and eventual disposal or decommissioning of a satellite \cite{CostEstimate}. Will affect users in how much we can charge
    \item \textbf{Durability:} satellite’s ability to withstand and continue functioning effectively under the harsh conditions of space for its intended mission lifespan \cite{DurabilityDef}.
    \item \textbf{Longevity:} satellite’s ability to remain operational and continue performing its intended functions over an extended period, ideally throughout or beyond its planned mission lifespan \cite{diebold2020longevity}.
    \item \textbf{Sustainability:} practices and technologies that minimize the environmental impact of satellite operations, both in space and on Earth. It also refers to the satellite's ability to self-sustain \cite{white2021sustainability}. 
    \item \textbf{Imaging:} the use of onboard sensors and cameras to capture images of space, planets, and celestial objects \cite{rees2013imaging}.
    \item \textbf{Communication:} refers to the process of transmitting and receiving data, signals, or messages between the satellite and ground stations \cite{Maral2011Communication}.
    \item \textbf{Power}: refers to the electrical energy generated, stored, and managed on-board to operate the satellite’s systems and payload. This further includes the power to launch the satellite and maintain it within its orbit \cite{Elbert2008Power}. 
\end{itemize}


\begin{table}[h]
\centering
\resizebox{\columnwidth}{!}{%
\begin{tabular}{lc|cc|cc|cc|}
\cline{3-8}
                                               & \multicolumn{1}{l|}{}                        & \multicolumn{2}{c|}{\textbf{ORACLE}}                                                                                 & \multicolumn{2}{c|}{\textbf{EVA}}                                                                                    & \multicolumn{2}{c|}{\textbf{NOVEA}}                                                                                  \\ \hline
\multicolumn{1}{|l|}{\textbf{Criteria}}         & \multicolumn{1}{l|}{\textbf{Weighted Rating}} & \multicolumn{1}{c|}{\textbf{Rating}} & \textbf{\begin{tabular}[c]{@{}c@{}}Weighted Rating x \\ Raiting\end{tabular}} & \multicolumn{1}{c|}{\textbf{Rating}} & \textbf{\begin{tabular}[c]{@{}c@{}}Weighted Rating x \\ Raiting\end{tabular}} & \multicolumn{1}{c|}{\textbf{Rating}} & \textbf{\begin{tabular}[c]{@{}c@{}}Weighted Rating x \\ Raiting\end{tabular}} \\ \hline
\multicolumn{1}{|l|}{\textbf{Imaging}}         & \textbf{7}                                   & \multicolumn{1}{c|}{3}               & 21                                                                            & \multicolumn{1}{c|}{3}               & 21                                                                            & \multicolumn{1}{c|}{5}               & 35                                                                            \\ \hline
\multicolumn{1}{|l|}{\textbf{Sustainability}} & \textbf{5}                                   & \multicolumn{1}{c|}{5}               & 25                                                                            & \multicolumn{1}{c|}{4}               & 20                                                                            & \multicolumn{1}{c|}{1}               & 5                                                                             \\ \hline
\multicolumn{1}{|l|}{\textbf{Longevity}}       & \textbf{4}                                   & \multicolumn{1}{c|}{5}               & 20                                                                            & \multicolumn{1}{c|}{3}               & 12                                                                            & \multicolumn{1}{c|}{2}               & 8                                                                             \\ \hline
\multicolumn{1}{|l|}{\textbf{Durability}}      & \textbf{3}                                   & \multicolumn{1}{c|}{3}               & 9                                                                             & \multicolumn{1}{c|}{3}               & 9                                                                             & \multicolumn{1}{c|}{5}               & 15                                                                            \\ \hline
\multicolumn{1}{|l|}{\textbf{Cost}}            & \textbf{6}                                   & \multicolumn{1}{c|}{4}               & 24                                                                            & \multicolumn{1}{c|}{3}               & 18                                                                            & \multicolumn{1}{c|}{2}               & 12                                                                            \\ \hline
\multicolumn{1}{|l|}{\textbf{Communication}}   & \textbf{1}                                   & \multicolumn{1}{c|}{5}               & 5                                                                             & \multicolumn{1}{c|}{5}               & 5                                                                             & \multicolumn{1}{c|}{3}               & 3                                                                             \\ \hline
\multicolumn{1}{|l|}{\textbf{Power}}           & \textbf{2}                                   & \multicolumn{1}{c|}{4}               & 8                                                                             & \multicolumn{1}{c|}{3}               & 6                                                                             & \multicolumn{1}{c|}{2}               & 4                                                                             \\ \hline
\multicolumn{2}{|l|}{{\ul \textbf{Total Value}}}                                              & \multicolumn{1}{c|}{29}              & 112                                                                           & \multicolumn{1}{c|}{24}              & 91                                                                            & \multicolumn{1}{c|}{20}              & 82                                                                            \\ \hline
\end{tabular}%
}
\caption{Decision Matrix}
\label{tab:my-table}
\end{table}
The decision as to which concept is chosen is listed above. Parameters derived from user needs are given ratings of importance and each concept is rated 1-5 on how well it succeeds in that parameter, 5 being the highest. We decided that the top 3 priorities for our users are imaging, cost, and sustainability. The rating is then multiplied with each weighted rating to give a score on each parameter for each concept. The sum of the products is then calculated to give us the best possible concept.\\
\vspace{1.5mm}

In our decision matrix, NOVEA was assigned a rating of 5 when it came to imaging as it is the most high resolution out of all three of the concepts since and due to its L2 orbit position in space it can take images without interruption as there is no eclipse. Also being in deep space it has a larger view of space and can source objects that are much further than that ORACLE or EVA can resolve. ORACLE achieved the highest sustainability rating due to its repairability due to its low earth orbit allowing for a modular design that can be changed, and repaired during its lifetime. The OCRACLE is also self-sustaining and uses onboard systems to produce its power. Its use of sustainable materials and practices makes it the clear option when it comes to sustainability. 
For similar reasons, the ORACLE achieved the highest rating in longevity due to material choice and its repairability. 

%Durability
Only NOVEA receives a rank of 5 due to its orbit location. Positioned at the L2 point, which is inaccessible for serviceability, NOVEA requires a higher safety factor for structural strength. Being further from Earth, it’s also more susceptible to collisions with undetected celestial debris due to its small size. To mitigate mission failure risks, the design should include thicker walls and panels to withstand potential impacts. ORACLE receives a rank of 3 based on their orbits. ORACLE is serviceable in LEO, but EVA has limited serviceability, restricted to a small window at perigee and the higher kinetic energy.


ORACLE was rated the highest for cost and power due to the relatively lower development cost in making a LEO satellite compared to a deep space satellite. It would also take less fuel to reach said orbit compared to the HEO or the L2 point of EVA and NOVEA respectively. The same can be said for communication for ORACLE and EVA as the orbit lies relatively close to the earth making communication arrays between the satellite and ground station more reliable compared to the NOVEA.

\subsection{Concluding Remarks}
The team has chosen ORACLE (First concept) as the optimal path forward. Its low risk, high redundancy, and proven record of Hubble offer favourable advantages over the other two options. To minimize risks and potential errors, the satellite shall adopt a similar orbital trajectory and operational framework to Hubble, with the key difference in components. The Hubble shall utilise current modern-era components to aid the mission, while also providing higher performance comparable to other current \& future satellites. The ORACLE additionally has a higher emphasis on sustainability, where the operations focus heavily on material and manufacturing practices. This approach leverages historical empirical data, utilizes components with a high technology readiness level, and capitalizes on Hubble’s existing infrastructure and framework. The goal is to develop a telescope that builds on Hubble’s legacy, smoothly integrating into the current system while enhancing performance and reliability compatible to the modern era.